\section{Constructing Proofs}

Although the rules are pretty simple, building proofs can get quite complicated.
Much like chess, the rules are easy to learn, but yon can get into trouble
quickly if you don't adopt some basic straegies to guide your moves.

\begin{observation}[Proof Strategies for TFL]
  \label{obs:tfl-proof-strategies}
  We can tersely describe the basic strategy for TFL proofs as follows.
  \begin{outline}[enumerate]

    \1 Work backwards from the conclusion.
    \2 Conjunction: what components needed to use the \eand I rule as the
    last step? i.e. what are the conjuncts?
    \2 Conditional: what are the assumption and conclusion of the subproof
    we need to take the last step? i.e. what are the antecedant and consequent
    of the conditional?
    \2 Negation: what subproof is needed to show a contradiction?
    \2 Disjunction: which of the disjuncts looks easiest to prove?
    
    \1 Work forwards from the current step.
    \2 Conjunction: unpack the conjuncts
    \2 Disjunction: what sentence do you need to prove from both disjuncts?
    \2 Conditional: how can you prove the antecedant?
    \2 Negation: how can you prove a contradiction?
    
    \1 Leverage \ebot
    \2 Work toward a contradiction to conclude anything, especially in a
    subproof
    \2 Indirect proof: build a contradictory subproof that assumes the negation

  \end{outline}
\end{observation}

Ideally, we will apply these stragegies in this order, wherever the rules are
appropriate to the current state of the proof. The \emph{work backwards} rules
in particular are intended as scaffolding to indicate which portions of the
overall proof need to be filled in. The rules themselves reflect this
partially-completed aspect.

\subsection{Working Backward}

Suffice it to say, taking account of what we are trying to prove is an important
component of successfully constructing proofs. This is the spirit in which the
rules for working backward are intended. They offer a road-map for proving the
conclusion.

Proving conjunctions amounts to proving each of the conjuncts.

\Needspace{7\baselineskip}
\begin{fitchrule}[Proving Conjunctions]
  \label{def:proving-conjunctions}
  \leavevmode
  \vspace{-0.9\baselineskip}
  \begin{fitchproof}
    \have[ ]     {1} {\vdots}    
    \have[m]     {2} {\ephi}
    \have[ ]     {3} {\vdots}    
    \have[n]     {4} {\epsi}   
    \have[n][+1] {5} {\ephi\eand\epsi} \ai{2,4}
  \end{fitchproof}
\end{fitchrule}

Proving conditionals and negations involves setting up subproofs of the
appropriate form.

\Needspace{5\baselineskip}
\twoitems{

\begin{fitchrule}[Proving Conditionals]
  \label{def:proving-conditionals}
  \leavevmode
  \vspace{-0.9\baselineskip}
  \begin{fitchproof}
    \open
    \hypo[m]     {1} {\ephi}          \by{AS}{}
    \have[ ]     {2} {\vdots}
    \have[n]     {3} {\epsi}
    \close
    \have[n][+1] {4} {\ephi\eif\epsi} \ii{1-3}
  \end{fitchproof}
\end{fitchrule}

}{
  
\begin{fitchrule}[Proving Negations]
  \label{def:proving-negations}
  \leavevmode
  \vspace{-0.9\baselineskip}
  \begin{fitchproof}
    \open
    \hypo[m]     {1} {\ephi}      \by{AS}{}
    \have[ ]     {2} {\vdots}
    \have[n]     {3} {\ebot}
    \close
    \have[n][+1] {4} {\enot\ephi} \ni{1-3}
  \end{fitchproof}  
\end{fitchrule}

}

Working backward from disjunctions is possible too, although the ambiguity
may lead to some backtracking. So this is a last resort.

\Needspace{4\baselineskip}
\begin{fitchrule}[Proving Disjunctions] There are two possible justifications
  for a disjunction, i.e. each of its disjuncts:\\
  \label{def:proving-disjunctions}
  \leavevmode
  %\vspace{-0.9\baselineskip}
  \twoitems{
    \begin{fitchproof}
      \have[ ]     {1} {\vdots}    
      \have[n]     {2} {\ephi}
      \have[n][+1] {3} {\ephi\eor\epsi} \oi{2}
    \end{fitchproof}
  }{
    \begin{fitchproof}
      \have[ ]     {1} {\vdots}    
      \have[n]     {2} {\epsi}
      \have[n][+1] {3} {\ephi\eor\epsi} \oi{2}
    \end{fitchproof}
  }
\end{fitchrule}

\subsection{Working Forward}

Working forward can be guided by the same schematic rules. For conjunction we
just unpack both of the conjuncts.

\begin{fitchrule}[Unpacking Conjunctions]
  \label{def:unpacking-conjunctions}
  \leavevmode
  \vspace{-0.9\baselineskip}
  \begin{fitchproof}
    \have[n]     {1} {\ephi\eand\epsi} 
    \have[ ][+1] {2} {\ephi}
    \have[n][+2] {3} {\epsi}
  \end{fitchproof}  
\end{fitchrule}

For disjunction we need to think more about what we are trying to do. Suppose
that \ephi and \epsi are the disjuncts and \echi is what we are trying to
prove (or a step in that direction). We will need a subproof from each of
the disjuncts to \echi.

\Needspace{10\baselineskip}
\begin{fitchrule}[Unpacking Disjunctions]
  \label{def:unpacking-disjunctions}
  \leavevmode
  \vspace{-0.9\baselineskip}
  \begin{fitchproof}
    \have[k]     {1} {\ephi\eor\epsi}
    \open
    \hypo[k][+1] {2} {\ephi} \by{AS}{}
    \have[ ]     {3} {\vdots}
    \have[m]     {4} {\echi}
    \close    
    \open
    \hypo[m][+1] {5} {\ephi} \by{AS}{}
    \have[ ]     {6} {\vdots}
    \have[n]     {7} {\echi}
    \close
    \have[n][+1] {8} {\echi} \oe{1,2-4,5-7}
  \end{fitchproof}  
\end{fitchrule}

Finally, conditionals and negations are pretty easy to unpack. But they each
require a proof (not necessarily via subproof) of a component (the antecedent
and the negated sentence, respectively).

\Needspace{6\baselineskip}
\twoitems{

\begin{fitchrule}[Unpacking Conditionals]
  \label{def:unpacking-conditionals}
  \leavevmode
  \vspace{-0.9\baselineskip}
  \begin{fitchproof}
    \have[m]     {1} {\ephi\eif\epsi} 
    \have[ ]     {2} {\vdots}
    \have[n]     {3} {\ephi}
    \have[n][+1] {4} {\epsi} \ie{1,3}
  \end{fitchproof}    
\end{fitchrule}

} {
  
\begin{fitchrule}[Unpacking Negations]
  \label{def:unpacking-negations}
  \leavevmode
  \vspace{-0.9\baselineskip}
  \begin{fitchproof}
    \have[m]     {1} {\enot\ephi} 
    \have[ ]     {2} {\vdots}
    \have[n]     {3} {\ephi}
    \have[n][+1] {4} {\ebot} \ne{1,3}
  \end{fitchproof}      
\end{fitchrule}

}

\subsection{Leveraging \ebot}

Sometimes the best way to get a subproof to its intended conclusion is to work
toward a contradiction. Then you can conclude the subproof with literally
anything, which is a powerfully flexible tool. Recall the proof of disjunctive
syllogism from the last lecture.

\begin{fitchproof}
  \hypo {1} {A \eor B}  \by{PR}{} 
  \hypo {2} {\enot A}   \by{PR}{}
  \open
  \hypo {3} {A}
  \have {4} {\ebot}     \ne{2,3}
  \have {5} {B}         \by{X}{4}
  \close
  \open
  \hypo {6} {B}
  \have {7} {B}         \by{R}{6}
  \close
  \have {7} {B}         \by{R}{1,3-5,6-7}
\end{fitchproof}      

Similarly, we can get any sentence we want in the main proof (or current
subproof) by assuming its negation and proving a contradiction. We will call
this technique \emph{indirect proof}. Its rule looks a lot like the identically
named rule of inference from the last lecture.

\begin{fitchrule}[Indirect Proof]
  \label{def:indirect-proof-rule}
  \leavevmode
  \vspace{-0.9\baselineskip}
  \begin{fitchproof}
    \open
    \hypo[m]     {1} {\enot\ephi} \by{AS}{}
    \have[ ]     {2} {\vdots}
    \have[n]     {3} {\ebot}
    \close
    \have[n][+1] {4} {\phi}       \by{IP}{1-3}
  \end{fitchproof}      
\end{fitchrule}

The method of indirect proof lets us prove all sorts of tautologies because it
depends on discharging the assumptions of a subproof. Aside from proving a
contradiction (explosion), this is the only way we can get ``something from
nothing'' in natural deduction.

\begin{fitchproof}
  \open
  \hypo {1} {\enot(A\eor\enot A)} \by{AS}{}
  \open
  \hypo {2} {A}                   \by{AS}{}
  \have {3} {A\eor\enot A}        \oi{2}
  \have {4} {\ebot}               \ne{1,3}
  \close
  \have {5} {\enot A}             \ni{2-4}
  \have {6} {A \eor\enot A}       \oi{5}
  \have {7} {\ebot}               \ne{1,6}
  \close
  \have {8} {A \eor\enot A}        \by{IP}{1-7}

\end{fitchproof}      

\subsection{Strategies in Action}

We will introduce some notation now, but postpone precise definitions until
Lecture 8 when we will focus our discussion on TFL metalogic. The concepts
should be more or less clear in their informal presentation.

\begin{enumerate}
\item When \epsi can be derived from \ephi, we will say \ephi\eproves\epsi
  (or \epsi\efollows\ephi) which is pronounced ``\ephi proves \epsi''.
\item We will use \eproves instead of \therefore, whenever appropriate.
  Specifically, we will use \ephi\eproves\epsi whenever \ephi\therefore\epsi is
  a (proof-theoretically) valid argument.
\item When \ephi\eproves\epsi and \ephi\efollows\epsi we can say that \ephi and
  \epsi are \emph{interderivable} and express this as
  \ephi\efollows\eproves\epsi.
\end{enumerate}

The only way to fully understand our proof strategies is to see them used on
some actual proofs. The interdreviability of the various formulas related to
material implication are a good choice, since that is an important property of
TFL and demonstrates most if not all of the strategies discussed earlier in
the lecture.

\begin{argument}
\item[ ] (A \eif B)\efollows\eproves\enot(A \eand\enot B)
\item[ ] (A \eif B)\efollows\eproves(\enot A \eor B)
\end{argument}

Proving the interderivability of two sentences typically amounts to proving both
directions of the derivation, one at a time.

\subsubsection{(A \eif B)\eproves \enot(A \eand\enot B)}

For (A \eif B)\eproves \enot(A \eand\enot B), we will work backward from a
negative. We just need to start with a template for a subproof of \ebot.

\begin{fitchproof}
  \hypo    {1} {A \eif B}             \by{PR}{}
  \open
  \hypo[i] {2} {A\eand\enot B}        \by{AS}{}
  \have[ ] {3} {\vdots}
  \have[j] {4} {\ebot}
  \close
  \have[n] {5} {\enot(A\eand\enot B)}  \ni{2-4}  
\end{fitchproof}      

We can start trying to get a contradiction by unpacking the conjunction in the
subproof. We'll fill that in.

\begin{fitchproof}
  \hypo    {1} {A \eif B}             \by{PR}{}
  \open
  \hypo[i] {2} {A\eand\enot B}        \by{AS}{}
  \have[ ] {3} {A}                    \ae{2}
  \have[ ] {4} {\enot B}              \ae{2}
  \have[ ] {5} {\vdots}
  \have[j] {6} {\ebot}
  \close
  \have[n] {} {\enot(A\eand\enot B)}  \ne{2-6}  
\end{fitchproof}      

We can fill in the missing lines using modus ponens and negation elimination.

\begin{fitchproof}
  \hypo {1} {A \eif B}             \by{PR}{}
  \open
  \hypo {2} {A\eand\enot B}        \by{AS}{}
  \have {3} {A}                    \ae{2}
  \have {4} {\enot B}              \ae{2}
  \have {5} {B}
  \have {6} {\ebot}                \ne{4,5}
  \close
  \have {7} {\enot(A\eand\enot B)}  \ni{2-6}  
\end{fitchproof}      

\subsubsection{\enot(A \eand\enot B)\eproves(A \eif B)}

For the backward case (\enot(A \eand\enot B)\eproves(A \eif B)), start by
setting up a proof schema for the conditional.

\begin{fitchproof}
  \hypo      {1} {\enot(A\eand\enot B)} \by{PR}{}
  \open
  \hypo[i]   {2} {A}                    \by{AS}{}
  \have[ ]   {3} { \vdots }
  \have[j]   {4} {B}
  \close
  \have[n]   {5} {A \eif B}             \ii{2-4}
\end{fitchproof}      

Now we just need to fill in the dots. Proving B from A will depend on how we
can use (1). Unpacking it and working forward will be a mess (a negated
conjunction is just a disjunction of negations). But perhaps we can proceed by
indirect proof, since (1) is a negation. We'll fill in some more with that
assumption in mind, opening a new subproof for IP.

\begin{fitchproof}
  \hypo      {1} {\enot(A\eand\enot B)} \by{PR}{}
  \open
  \hypo      {2} {A}                    \by{AS}{}
  \open
  \hypo      {3} {\enot B}              \by{AS}{}
  \have      {4} {\vdots}
  \have[i]   {5} {\ebot}
  \close
  \have[j]   {6} {B}
  \close
  \have[n]   {7} {A \eif B}             \ii{2-6}
\end{fitchproof}      

And we're done if we just use the \eand introduction rule to fill in the missing
line.

\begin{fitchproof}
  \hypo {1} {\enot(A\eand\enot B)} \by{PR}{}
  \open
  \hypo {2} {A}                    \by{AS}{}
  \open
  \hypo {3} {\enot B}              \by{AS}{}
  \have {4} {A \eand\enot B}       \ai{2,3}
  \have {5} {\ebot}                \ne{1,4}
  \close
  \have {6} {B}                    \by{IP}{3-5}
  \close
  \have {7} {A \eif B}             \ii{2-6}
\end{fitchproof}      

\subsubsection{(A \eif B)\eproves(\enot A \eor B)}

Another well-known equivalence is (A \eif B)\emutual(\enot A \eor B). Start with
(A \eif B)\eproves(\enot A \eor B). To work backward from a conclusion would
require choosing a disjunct, so might be a little open-ended. We'll get back to
that if we need to. Let's try indirect proof instead.

\begin{fitchproof}
  \hypo    {1} {A \eif B}              \by{PR}{}
  \open
  \hypo[i] {2} {\enot(\enot A \eor B)} \by{AS}{}
  \have[ ] {3} {\vdots}
  \have[j] {4} {\enot A \eor B}
  \have[k] {5} {\ebot}                 \ne{1,4}
  \close
  \have[n] {6} {\enot A \eor B}        \by{IP}{2-5}
\end{fitchproof}      

To get a contradiction we want to prove either \enot A or B to we can say
\enot A\eor B. If we assume A we can use use modus ponens to unpack (1), so
we'll go that route (i.e. we'll prove \enot A). First sketch out the subproof.

\begin{fitchproof}
  \hypo    {1} {A \eif B}              \by{PR}{}
  \open
  \hypo    {2} {\enot(\enot A \eor B)} \by{AS}{}
  \open
  \hypo[i] {3} {A}                     \by{AS}{}
  \have[ ] {4} {\vdots}
  \have[j] {5} {\ebot}                 
  \close
  \have[k] {6} {\enot A}               \ni{3-5}
  \have[l] {7} {\enot A \eor B}        \oi{6}
  \have[m] {8} {\ebot}                 \ne{2,7}
  \close
  \have[n] {9} {\enot A \eor B}        \by{IP}{2-8}
\end{fitchproof}      

Then fill it in by applying \eif elimination to the premise using the new
assumption.

\begin{fitchproof}
  \hypo {1} {A \eif B}              \by{PR}{}
  \open
  \hypo {2} {\enot(\enot A \eor B)} \by{AS}{}
  \open
  \hypo {3} {A}                     \by{AS}{}
  \have {4} {B}                     \ie{1,3}
  \have {5} {\enot A \eor B}        \oi{4}
  \have {6} {\ebot}                 \ne{2,5}
  \close
  \have {7} {\enot A}               \ni{3-6}
  \have {8} {\enot A \eor B}        \oi{7}
  \have {9} {\ebot}                 \ne{2,8}
  \close
  \have {10}{\enot A \eor B}        \by{IP}{2-9}
\end{fitchproof}      

\subsubsection{(\enot A \eor B)\eproves(A \eif B)}

Now the other direction: (\enot A \eor B)\eproves(A \eif B). So far we haven't
seen explosion or \eor unpacked, but this proof has them both. First, we note
that we are trying to prove A \eif B by adding a template for the appropriate
subproof.

\begin{fitchproof}
  \hypo    {1} {\enot A \eor B} \by{PR}{}
  \open
  \hypo[i] {2} {A}              \by{AS}{}
  \have[ ] {3} {\vdots}
  \have[j] {4} {B}
  \close 
  \have[n] {5} {(A \eif B)}     \ii{2-4}
\end{fitchproof}      

The premise is a disjunction, so if it's gonna play a role in the proof, we're
gonna need to unpack it with \eor elimination, so let's set up the subproofs for
that now.

\begin{fitchproof}
  \hypo    {1} {\enot A \eor B} \by{PR}{}
  \open
  \hypo    {2} {A}              \by{AS}{}
  \open
  \hypo[i] {3} {\not A}         \by{AS}{}
  \have[ ] {4} {\vdots}
  \have[j] {5} {B}
  \close
  \open
  \hypo[k] {6} {B} 
  \have[ ] {7} {\vdots}
  \have[l] {8} {B}
  \close 
  \have[m] {9} {B}              \oe{1,3-5,6-8}
  \close  
  \have[n] {10}{(A \eif B)}     \ii{2-9}
\end{fitchproof}      

Both subproofs follow immediately. The first is a contradiction, so we can apply
explosion to get B. The second is just repetition; in fact we don't even need
the placeholder line in the template.

\begin{fitchproof}
  \hypo {1} {\enot A \eor B} \by{PR}{}
  \open
  \hypo {2} {A}              \by{AS}{}  
  \open
  \hypo {3} {\not A}         \by{AS}{}
  \have {4} {\ebot}          \ne{2,3}
  \have {5} {B}              \by{X}{4}
  \close
  \open
  \hypo {6} {B}              \by{AS}{}
  \have {7} {B}              \by{R}{6}
  \close  
  \have {8} {B}              \oe{1,3-5}
  \close
  \have {9}{(A \eif B)}      \ii{2-8}
\end{fitchproof}      

